\subsection{The Radical of a Tensor Product of Algebras}

Let \(A_{1}\) and \(A_{2}\) be K-algebras, and let \(A = A_{1} \otimes A_{2}\) .

PROPOSITION 3. --- Suppose that the algebras \(\mathrm{A}_1\) and \(\mathrm{A}_2\) are semisimple, with respective centers \(\mathbf{Z}_1\) and \(\mathbf{Z}_2\) . Set \(\mathbf{Z} = \mathbf{Z}_1 \otimes \mathbf{Z}_2\) .

\begin{enumerate}
\def\labelenumi{\alph{enumi})}
\item
  The mapping \(a \mapsto aA\) is an isomorphism from the set of ideals of Z, ordered by inclusion, to the set of two-sided ideals of A, ordered by inclusion.
\item
  The radical of A is equal to the intersection of the maximal two-sided ideals of A and equals \(\Re(Z)A\) .
\item
  If one of the K-algebras \(Z_{1}\) and \(Z_{2}\) is separable, in particular if the field K is perfect, then the radicals of the rings Z and A are reduced to 0.
\end{enumerate}

Each algebra \(A_{i}\) is the product of finitely many simple algebras. Now, the center of a product of rings is the product of the centers, and we have the analogous assertions for radicals (VIII, p.~156, Corollary 3) and for two-sided ideals (I, §8, No.~10, p.~109, Proposition 8). It therefore suffices to prove Proposition 3 under the assumption that \(A_{1}\) and \(A_{2}\) are simple algebras.

For \(i \in \{1,2\}\) , set \(B_i = A_i \otimes A_i^o\) , and view \(A_i\) as a \(B_i\) -module, where the homotheties are given by the formula

\[
(x \otimes y) z = x z y
\]

for x, y, and z in \(A_{i}\) . The commutant of the \(B_{i}\) -module \(A_{i}\) is \((\mathrm{Z}_{i})_{\mathrm{A}_{i}}\) , the set of homotheties by elements of \(Z_{i}\) ; we identify it with \(Z_{i}\) . Moreover, the \(B_{i}\) -submodules of \(A_{i}\) are the two-sided ideals of \(A_{i}\) , and since the ring \(A_{i}\) is simple, the only two-sided ideals are 0 and \(A_{i}\) . Hence \(A_{i}\) is a simple \(B_{i}\) -module. Moreover, the \((\mathrm{B}_{1}\otimes\mathrm{B}_{2})\) -submodules of \(A_{1}\otimes A_{2}\) are the two-sided ideals of the ring \(A_{1}\otimes A_{2}\) .

Assertion a) therefore follows from VIII, p.~214, Theorem 2, b) applied to the simple \(\mathrm{B}_1\) -module \(\mathrm{A}_1\) , with commutant \(\mathrm{Z}_1\) , and to the simple \(\mathrm{B}_2\) -module \(\mathrm{A}_2\) , with commutant \(\mathrm{Z}_2\) .

Let us prove assertion b). The intersection of the maximal two-sided ideals of A is the radical of the \((\mathrm{B}_{1}\otimes\mathrm{B}_{2})\) -module \(A_{1}\otimes A_{2}\) ; by Corollary 2 of VIII, p.~215, this intersection coincides with \(\Re(\mathrm{Z})\mathrm{A}\) . The algebras \(Z_{1}\) and \(Z_{2}\) are commutative and semisimple. The radical of the ring Z therefore consists of nilpotent elements (VIII, p.~215, Theorem 3), and the two-sided ideal \(\Re(\mathrm{Z})\mathrm{A}\) of the ring A is contained in the radical \(\Re(\mathrm{A})\) (VIII, p.~157, Remark 1). However, the intersection of the maximal two-sided ideals of A contains \(\Re(\mathrm{A})\) (VIII, p.~155, Proposition 5, d)). This proves assertion b).

The tensor product of a separable commutative algebra and a reduced commutative algebra is a reduced ring (V, §15, No.~2, p.~120, Proposition 5). The algebras \(Z_1\) and \(Z_2\) are commutative and semisimple, hence reduced. If one of the algebras \(Z_1\) and \(Z_2\) is separable, then the algebra \(Z\) is reduced; it is therefore without radical by Theorem 3 of VIII, p.~215, and we have \(\Re(A) = \Re(Z)A = 0\) . This is, in particular, the case when the field \(K\) is perfect because every reduced commutative algebra over a perfect field is separable (V, §15, No.~5, p.~125, Theorem 3).

COROLLARY. --- Suppose that the algebras \(A_{1}\) and \(A_{2}\) are simple and that the center \(Z_{1}\) of \(A_{1}\) is reduced to K. Then the ring \(A_{1} \otimes A_{2}\) has no other two-sided ideals than 0 and itself.

By assumption, we have \(Z_{1} = K\) , and since \(A_{2}\) is simple, its center \(Z_{2}\) is a field (VIII, p.~121, Corollary 1, a)). The ring \(Z = Z_{1} \otimes Z_{2}\) is therefore a field, and the corollary follows from Proposition 3, a).

